\section{The master space and its fixed theories}\label{geo:section}

We place the base and its projectivization inside a single smooth
projective variety with a torus action.  Its fixed components will carry
inverse-Euler twists of their ordinary theories.  This section specifies
their curve labels and pairings, and recalls the ancestor and fundamental
solution conventions needed for the localization comparison.

\subsection{Geometry and integral curve labels}

Tensoring $E$ by a line bundle does not change the projective bundle of
one-dimensional subspaces.  Choose such a twist so that $E^*$ is globally
generated, and continue to denote the resulting bundle by $E$.
Set
\[
       X=\mathbb P_B(E),\qquad W=\mathbb P_B(E\oplus\mathcal O_B).
\]
Let $T=\C^*$ act with weight zero on $E$ and weight one on
$\mathcal O_B$, and let $\lambda\in H_T^2(\mathrm{pt})$ be the class of
the weight-one representation.  The fixed locus is
$W^T=X\sqcup B$, where the second component is the section corresponding
to the summand $\mathcal O_B$.
Write $\pi$ for either projective-bundle projection and put
$p=c_1^T(\mathcal O_W(1))$.  The ordinary normal bundles and their torus
weights are
\begin{equation}
\begin{array}{c|c|c|c|c|c}
 F&N_F&n_F=\rk N_F&w_F&W_F=n_Fw_F&j_F\\ \hline
 X&\mathcal O_X(1)&1&1&1&0\\
 B&E&r&-1&-r&1
\end{array}
\label{geo:normal-data}
\end{equation}
The restrictions of the divisor are
\[
              p|_X=c_1(\mathcal O_X(1)),\qquad p|_B=-\lambda.
\]
Here $w_F$ is the common torus weight on the normal fibers; the integers
$j_F$ are recorded for the fixed-section curve factors in the shift
operator below.  The normal-bundle descriptions follow from the relative
tangent space $\Hom(L,V/L)$ at a line $L\subset V$.

\begin{lemma}\label{geo:curve-labels}
For $P=X$ or $W$, the map
\begin{equation}
 H_2(P;\Z)\longrightarrow H_2(B;\Z)\oplus\Z,
 \qquad d\longmapsto
       \left(\pi_*d,\int_d c_1(\mathcal O_P(1))\right)
 \label{geo:integral-splitting}
\end{equation}
is an isomorphism, including torsion.  An effective class has an effective
base pushforward $\beta$ and a nonnegative tautological degree $l$.
\end{lemma}

\begin{proof}
The structure group of a projective bundle acts trivially on the integral
homology of its fiber.  In total degree two the homology Serre spectral
sequence therefore has only the base term $H_2(B;\Z)$ and the fiber term
$H_2(\mathbb P^{s-1};\Z)=\Z$, where $s\geq2$ is the bundle rank.
The integral class $c_1(\mathcal O_P(1))$ pairs to one with the fiber line.
Consequently the fiber generator survives: its image in total homology
cannot be zero or a nontrivial multiple quotient, as either possibility
would contradict that pairing.  The edge filtration gives an exact
sequence
\[
 0\longrightarrow\Z[\text{fiber line}]
   \longrightarrow H_2(P;\Z)
   \xrightarrow{\pi_*}H_2(B;\Z)\longrightarrow0.
\]
Pairing with $c_1(\mathcal O_P(1))$ is a retraction on its first term;
together with $\pi_*$ it gives \eqref{geo:integral-splitting}.
This integral argument does not discard torsion.
Finally, the dual bundles $E^*$ and $E^*\oplus\mathcal O_B$ are globally
generated, so the corresponding tautological line bundles are globally
generated.  Their degrees on effective curves are nonnegative, and
proper pushforward preserves effective curve classes.
\end{proof}

We write $Q^\beta y^l$ for the actual class corresponding to $(\beta,l)$.
Under the inclusion $X\hookrightarrow W$ these two coordinates are
unchanged, and under $B\hookrightarrow W$ a class $\beta$ becomes
$(\beta,0)$.  Thus the fixed theories and the master-space theory have
compatible full integral labels.  We may allow all pairs of an effective
base class and an integer $l\geq0$ as formal labels, assigning coefficient
zero when a pair is not effective for the target in question.

Choose an ample integral class $H$ on $B$.  For a sufficiently large
integer $M$, the class $c_1(\mathcal O_P(1))+M\pi^*H$ is ample for both
projective bundles.  We complete in the degree
\[
                         l+M\int_\beta H.
\]
There are finitely many effective labels of bounded degree, as in
Section~\ref{conv:section}.  Unless specified otherwise, positive
Novikov filtration means positive degree for this completion, so it
includes positive fiber degree even when the base class is zero.
The separate completion by positive base degree will be used only when
the fiber parameter $y$ is treated as a variable on a complex domain.

\subsection{Equivariant pairings and inverse-Euler twists}

The equivariant projective-bundle formula makes $H_T^*(W)$ free over
$\C[\lambda]$, with basis
\[
                    \pi^*\phi_a\,p^j,\qquad 0\leq j\leq r,
\]
for a homogeneous Hodge basis $(\phi_a)$ of $H_B$.
Its equivariant integration pairing is perfect over $\C[\lambda]$.
Indeed, integration over the projective fibers gives an antitriangular
matrix in the powers of $p$, whose antidiagonal blocks are the ordinary
Poincar\'e pairing of $B$.

After extending scalars to $\C(\lambda)$, restriction to the fixed locus
is an isomorphism of paired spaces
\begin{equation}
\begin{gathered}
 \left(H_T^*(W),\eta_W\right)\otimes_{\C[\lambda]}\C(\lambda)
   \cong
 \bigoplus_{F=B,X}\left(H_F\otimes\C(\lambda),\eta_F^t\right),\\
 \eta_F^t(a,b)=\int_F\frac{a\cup b}{e_T(N_F)}.
\end{gathered}
 \label{geo:fixed-pairing}
\end{equation}
The unit restricts to the sum of the two fixed units.  Every normal weight
is nonzero, so the Euler classes in this formula are invertible over
$\C(\lambda)$.

The corresponding twisted Gromov--Witten theory on $F$ inserts
\begin{equation}
    e_T\!\left(R^\bullet\rho_*f^*N_F\right)^{-1}
    \label{geo:Euler-twist}
\end{equation}
in each stable-map integral.  Here $\rho$ and $f$ are the universal curve
and universal map, and the inverse Euler class is extended
multiplicatively to the indicated $K$-theory class.  The torus acts
trivially on $F$ and with weight $w_F$ on $N_F$.
We use subscripts $F,\mathrm{tw}$ for its potentials and fundamental
solutions.  Its degree-zero three-point pairing is exactly $\eta_F^t$.

Give $p$ and $\lambda$ Hodge bidegree $(1,1)$.  In a homogeneous polynomial
basis the first-Hodge grading acts on basis vectors, while differentiation
in $\lambda$ records the degree of equivariant coefficients when needed.
The virtual dimension identities can be expressed in this first degree:
algebraicity forces the sum of the Hodge-degree differences of insertions
in a nonzero invariant to be zero.  This remains true equivariantly, by
finite-dimensional approximations or by fixed-locus integration.
It is distinct from the ordinary cohomological-degree count used to bound
fiber degree later.

\subsection{Ancestors and the fundamental solution}\label{geo:ancestors}

We give the same definitions for an ordinary theory, the equivariant
theory of $W$, or one of the twisted fixed theories.  Denote its pairing
by $\eta$, its Novikov monomial for a class $d$ by $\mathsf q^d$, and
its stable-map integrals by brackets.  Let $u$ be a formal even primary
background.  For distinguished insertions $a_1,\ldots,a_m$, set
\[
 \langle a_1,\ldots,a_m\rangle_{g,u}
   =\sum_{\substack{d\ \mathrm{effective}\\n\geq0}}
      \frac{\mathsf q^d}{n!}
        \langle a_1,\ldots,a_m,
                   \underbrace{u,\ldots,u}_{n}\rangle_{g,m+n,d},
\]
where $d$ runs over effective classes and only stable-map terms are
included.  Descendant insertions in these brackets use the map cotangent
classes.

For $2g-2+m>0$, there is a stabilization map
\[
 \operatorname{st}_{g,m}:
 \overline{\mathcal M}_{g,m+n}(Y,d)
     \longrightarrow\overline{\mathcal M}_{g,m}
\]
which forgets the map and the $n$ background marks and stabilizes the
remaining curve.  Define $\bar\psi_i=\operatorname{st}_{g,m}^*\psi_i$ for
$1\leq i\leq m$.  The ancestor potential at background $u$ is
\begin{equation}
\begin{split}
 \overline F_g(u;t)
  &=\sum_{\substack{d\ \mathrm{effective},\ m,n\geq0\\2g-2+m>0}}
     \frac{\mathsf q^d}{m!n!}
       \left\langle
        \prod_{i=1}^m\left(\sum_{j\geq0}\bar\psi_i^j\ev_i^*t_j\right)
        \prod_{a=m+1}^{m+n}\ev_a^*u
       \right\rangle_{g,m+n,d},\\
 \mathcal A(u;t)&=
       \exp\left(\sum_{g\geq0}\hbar^{g-1}\overline F_g(u;t)\right).
\end{split}
\label{geo:ancestor-potential}
\end{equation}
The bracket in this display denotes integration of the displayed product,
with the twist when appropriate.  All terms with curve-unstable
distinguished data are omitted, even when the stable-map space itself
exists.  The genus-one term with no distinguished marks is therefore
omitted as well.  Since the ancestors come from
$\overline{\mathcal M}_{g,m}$, their total power exceeds the dimension
$3g-3+m$ only when the corresponding product vanishes.

\begin{definition}\label{geo:fundamental}
The fundamental solution $S(u,z)=\id+O(z^{-1})$ is defined by
\begin{equation}
 \eta(b,S(u,z)a)
    =\eta(b,a)+\left\langle b,\frac{a}{z-\psi}\right\rangle_{0,u},
 \qquad
 \frac1{z-\psi}=\sum_{j\geq0}\psi^jz^{-j-1}.
 \label{geo:S-definition}
\end{equation}
The pairing term supplies the unstable two-point contribution.
\end{definition}

The genus-zero splitting and cotangent-class identities give
\begin{equation}
 S(u,-z)^*S(u,z)=\id,
 \qquad z\partial_v S(u,z)=(v\star_u)S(u,z),
 \label{geo:S-identities}
\end{equation}
where $\star_u$ is the quantum product and $\partial_v$ differentiates
the primary background.  For clarity, the genus-zero cone $\mathcal L$
is the formal Lagrangian graph of $dF_0$ in the Darboux coordinates
$(q,p)$, with $q=t-z1$.  Its point with descendant input $\epsilon$ has
positive coordinate $\epsilon-z1$, and the tangent space there is the
graph of the Hessian of $F_0$ at $\epsilon$.  The genus-zero string,
dilaton, and topological
recursion identities say that each tangent space $T$ is tangent along
$zT\subset\mathcal L$.  At primary coordinate $u$ that tangent space is
$S(u,z)^{-1}\cH_+$; see \cite{GiventalCone}.
The ancestor--descendant correspondence in our group-action convention is
\begin{equation}
                    \mathcal A(u)=c(u)\,U_{S(u)}Z,
                    \label{geo:ancestor-descendant}
\end{equation}
where $c(u)$ is independent of ancestor variables.  These formulas follow
from the splitting identity for $\psi_i-\bar\psi_i$ and its genus-zero
specializations; see \cite{KontsevichManin,Getzler,GiventalQuant}.
They hold with the super contractions of Section~\ref{conv:section} and
with the inverse-Euler twists above.

We use \eqref{geo:ancestor-descendant} only on the fixed theories, with
$u$ of positive total Novikov filtration.  There $S(u,z)-\id$ has positive
filtration, and every Novikov coefficient has only finitely many negative
powers of $z$: only finitely many background marks occur at that
coefficient, and the map cotangent classes on the fixed targets are
nilpotent.  Thus the quantized action is covered by
Lemma~\ref{conv:finite-support}.
For the master space we abbreviate
\[
               S_W=S_W(0,z),\qquad\mathcal A_W=\mathcal A_W(0).
\]
Its rational loop-variable expansions and the quantized action used to
compare it with the fixed theories will be established directly by
localization in the next section.
